\documentclass[10pt,a4paper]{amsart}
\usepackage[margin=19mm]{geometry}
\usepackage{amsmath,amssymb,mathtools}
\usepackage[scr=rsfs,cal=euler]{mathalfa}
\usepackage{xfrac}
\usepackage[unicode,hidelinks,bookmarks=false]{hyperref}
\newcommand{\ie}{i.e.\ }
\newcommand{\cf}{cf.\ }
\newcommand{\resp}{respectively\ }
\newcommand{\Subsection}{\subsection}
\providecommand{\nobreakdash}{\nobreak\hbox{-}}
\setlength{\emergencystretch}{2em}
\multlinegap0pt
\usepackage{dsfont}
\usepackage{amssymb}
\newtheorem{lemma}{Lemma}
\newtheorem{proposition}{Proposition}
\newtheorem{theorem}{Theorem}
\DeclareMathOperator{\supp}{supp}
\title{Hierarchical symmetry selects log-Poisson cascades: classification, uniqueness, and stability}
\author{E. M. Freeburg}
\date{Source: arXiv:2604.01632v2 (CC BY 4.0)}
\begin{document}
\maketitle

\noindent\emph{Source and licence.} Hierarchical symmetry selects log-Poisson cascades: classification, uniqueness, and stability. Version arXiv:2604.01632v2 (CC BY 4.0), distributed by arXiv under the Creative Commons Attribution 4.0 International licence, http://creativecommons.org/licenses/by/4.0/, as recorded in the arXiv licence metadata for this exact version and shown on the abstract page https://arxiv.org/abs/2604.01632v2. Full licence notice: \texttt{licenses/arxiv-2604.01632-CC-BY-4.0.txt}.\par\smallskip

\noindent\textit{Editorial scope.} Counted unit: the article's theorem thm:classification (source0.tex lines 612--624). The article's complete original proof of it is delivered with it (source0.tex lines 626--758, one \texttt{proof} environment of 133 lines). No mathematics is written, repaired or re-derived: the statement and the proof are the article's own lines, in the article's own notation, and no retained line is substituted or re-flowed.  Retained uncounted as the source-local context the proof needs: lines 282--286; lines 325--338; lines 562--567.  Nothing was omitted from any retained span, so the package carries no omission record.  No retained line contains a citation, so the wrapper carries no bibliography and no bibliography entry.  No retained line contains a figure environment, an image-inclusion command or drawing code, so no image is delivered and the package declares no asset dependency.  Editorial formatting only, in addition: an AMS \texttt{amsart} preamble that restates the article's own declarations and macros used by the retained lines, namely the environments \texttt{lemma}, \texttt{proposition}, \texttt{theorem} on separate counters, together with the standard packages those lines need.  The retained environments print in this document as Lemma 1, Proposition 1, Theorem 1 in order of appearance, whereas the article prints the same items with the numbering of its own full document.  The complete native source of the article as distributed in the arXiv e-print for this exact version is bundled unchanged as \texttt{sources/197678/source0.tex}.
\begin{align}\label{eq:A1}
  \delta_{p+k}
  = (1 - \beta)\,L + \beta\,\delta_p.
  \tag{$*$}
\end{align}

\begin{lemma}[Exponent Form]\label{lem:exponent}
If the incremental exponents $\{\delta_p\}$ satisfy the A1
recurrence~\eqref{eq:A1} with $\beta \in (0,1)$, then with $\zeta_0
= 0$:
\begin{align}\label{eq:zeta}
  \zeta_p = \gamma p + C\bigl(1 - \beta^{p/k}\bigr),
  \tag{$**$}
\end{align}
where $\gamma = \delta_\infty / k$ and $C = (\delta_0 -
\delta_\infty)/(1 - \beta)$.

We emphasize that this lemma is purely algebraic and involves no
probabilistic content.
\end{lemma}

\begin{proposition}[Converse]\label{prop:converse}
If the cascade multiplier $W$ is log-Poisson---\allowbreak that is,
$\log W = a + bN$ with $N \sim \mathrm{Poisson}(\lambda)$, $b < 0$, $\lambda >
0$---then the incremental scaling exponents satisfy A1 with $\beta =
e^{bk} \in (0,1)$.
\end{proposition}

\begin{theorem}[Log-ID Classification]\label{thm:classification}
Let $\{W_n\}$ be an i.i.d.\ multiplicative cascade with nontrivial
intermittency ($C > 0$), whose generator $\log W$ is infinitely
divisible with L\'evy triplet $(a, \sigma^2, \nu)$.  Then A1 holds
with $\beta \in (0,1)$ if and only if $\sigma^2 = 0$ and $\nu =
\lambda\delta_b$ for some $b < 0$, $\lambda > 0$.  That is:
\begin{center}
\emph{A1 selects exactly the log-Poisson class from the full
log-infinitely-divisible family.}
\end{center}
No other log-ID cascade---log-normal, log-stable, or any
intermediate---satisfies A1.
\end{theorem}

\begin{proof}
\textit{Reverse direction.} If $\nu = \lambda\delta_b$ with $b < 0$
and $\sigma^2 = 0$, then $\log W = a + bN$ with $N \sim
\mathrm{Poisson}(\lambda)$, and A1 holds by
Proposition~\ref{prop:converse}.

\textit{Forward direction.}  Assume A1 holds.  We show $\sigma^2 = 0$
and $\nu = \lambda\delta_b$.

\textit{Step~1 (unsplit form).}  The cumulant generating function of
$\log W$ is
\[
  \psi(p) = ap + \frac{\sigma^2 p^2}{2}
  + \int\bigl(e^{px} - 1 - px\,\mathds{1}_{|x|\leq 1}\bigr)
  \,\nu(dx),
\]
finite for all $p \ge 0$ since all moments of $W$ are finite.  With
$\zeta_p = \psi(p)/\ln r$ and $\delta_p = (\psi(p+k) -
\psi(p))/\ln r$, define $\phi(p) = \psi(p+k) - \psi(p)$.  Then
\[
\begin{aligned}
  \phi(p) &= ak + \sigma^2 k\!\left(p + \tfrac{k}{2}\right)
  + \int g_p(x)\,\nu(dx),\\
  g_p(x) &:= e^{px}\bigl(e^{kx}-1\bigr) - kx\,\mathds{1}_{|x|\leq 1},
\end{aligned}
\]
where $g_p$ is $\nu$-integrable for each $p$, being the difference of
the two compensated L\'evy--Khintchine integrands.  \emph{No
splitting of the integral is performed at this stage.}

\textit{Step~1$'$ (sign inventory).}  For every $p \ge 0$:
\begin{itemize}
\item on $(0,1]$: $e^{px} \ge 1$ gives $g_p(x) \ge (e^{kx}-1) - kx
  \ge 0$, and $g_p(x) \uparrow \infty$ pointwise as $p \to \infty$;
\item on $(1,\infty)$: $g_p(x) = e^{px}(e^{kx}-1) \ge 0$, increasing
  to $+\infty$ pointwise;
\item on $[-1,0)$: $|e^{px}(e^{kx}-1)| \le 1 - e^{kx} \le k|x|$, so
  $0 \le g_p(x) \le k|x|$, with $g_p(x) \to k|x|$ pointwise as
  $p \to \infty$;
\item on $(-\infty,-1)$: $-1 \le g_p(x) \le 0$, with $g_p(x) \to 0$
  pointwise.
\end{itemize}
In particular $\int g_p\,d\nu \ge -\nu((-\infty,-1))$, uniformly in
$p$.

\textit{Step~2} ($\sigma^2 = 0$).  If $\sigma^2 > 0$ then, by
Step~1$'$,
\[
  \phi(p) \;\ge\; ak + \sigma^2 k\Bigl(p + \tfrac{k}{2}\Bigr)
  - \nu\bigl((-\infty,-1)\bigr) \;\longrightarrow\; +\infty,
\]
so $\delta_p = \phi(p)/\ln r \to -\infty$, contradicting the finite
limit $\delta_\infty$ forced by A1.  Hence $\sigma^2 = 0$.
\textit{This eliminates all log-normal and mixed Gaussian-jump
generators.}

\textit{Step~3} ($\supp(\nu) \subseteq (-\infty,0]$).  If $\nu$ has
mass on $(0,\infty)$ then, since $g_p \ge 0$ there and $g_p \uparrow
\infty$ pointwise, monotone convergence gives $\int_{(0,\infty)}
g_p\,d\nu \to \infty$, while $\int_{(-\infty,0)} g_p\,d\nu \ge
-\nu((-\infty,-1))$; again $\phi(p) \to \infty$ and $\delta_p \to
-\infty$, a contradiction.  Therefore $\supp(\nu) \subseteq
(-\infty,0]$.  \textit{This eliminates all generators with positive
jumps.}

\textit{Step~3\textonehalf} (integrability near~$0$).  We claim A1
forces $\int_{[-1,0)} |x|\,\nu(dx) < \infty$.  With $\sigma^2 = 0$
and $\supp\nu \subseteq (-\infty,0]$, apply Fatou's lemma on
$[-1,0)$ (integrand $g_p \ge 0$ by Step~1$'$, pointwise limit
$k|x|$) and dominated convergence on $(-\infty,-1)$ (bounded by~$1$,
finite mass):
\[
  \liminf_{p\to\infty} \phi(p)
  \;\ge\; ak + k\int_{[-1,0)} |x|\,\nu(dx)
  - \nu\bigl((-\infty,-1)\bigr).
\]
If $\int_{[-1,0)}|x|\,d\nu = \infty$ then $\phi(p) \to \infty$ and
$\delta_p \to -\infty$, contradicting A1.  Hence
$\int_{|x|\le1}|x|\,d\nu < \infty$, the compensator integral
$k\int x\,\mathds{1}_{|x|\le1}\,d\nu$ is finite, and \emph{only now}
may the integral be split:
\[
  \phi(p) = c_0 + \int_{(-\infty,0)} e^{px}\bigl(e^{kx}-1\bigr)
  \,\nu(dx),
  \qquad
  c_0 = ak - k\!\int x\,\mathds{1}_{|x|\leq 1}\,\nu(dx).
\]
The remaining integrand is dominated by $1 - e^{kx} \le
\min(1, k|x|) \in L^1(\nu)$ and tends to $0$ pointwise, so by
dominated convergence $\phi(p) \to c_0$ as $p \to \infty$; thus
$\phi_\infty = c_0$ and $\delta_\infty = c_0/\ln r$.

\textit{Step~4} ($\nu$ is a single Dirac mass).  A1 at $p = mk$
gives, by Lemma~\ref{lem:exponent}(2),
$\delta_{mk} - \delta_\infty = (\delta_0-\delta_\infty)\beta^m$;
multiplying by $\ln r$,
\begin{equation}\label{eq:dagger}
  \int_{(-\infty,0)} e^{mkx}\bigl(e^{kx}-1\bigr)\,\nu(dx)
  = A\beta^m
  \qquad\text{for all } m \geq 0,
\end{equation}
where $A = (\delta_0 - \delta_\infty)\ln r < 0$ (nontrivial
intermittency and $\ln r < 0$).  Substitute $u = e^{kx}$, mapping
$(-\infty,0) \to (0,1)$; let $\tilde\nu$ be the pushforward of $\nu$
and set
\[
  \eta := (1-u)\,d\tilde\nu \;\ge\; 0,
  \qquad
  \mu_m := \int_{(0,1)} u^m \, d\eta .
\]
Then \eqref{eq:dagger} reads $\mu_m = |A|\,\beta^m$ for all $m \ge
0$; the case $m=0$ shows $\eta$ is a finite positive measure of
total mass $|A|$.  Only $m \in \{0,1,2\}$ are needed:
\[
  \int_{(0,1)} (u-\beta)^2 \, d\eta
  \;=\; \mu_2 - 2\beta\mu_1 + \beta^2\mu_0
  \;=\; |A|\bigl(\beta^2 - 2\beta^2 + \beta^2\bigr)
  \;=\; 0 .
\]
Since $(u-\beta)^2 > 0$ on $(0,1)\setminus\{\beta\}$ and $\eta \ge
0$, we conclude $\eta\bigl((0,1)\setminus\{\beta\}\bigr) = 0$ and
$\eta(\{\beta\}) = |A|$.  Because $u - 1 \neq 0$ on $(0,1)$, the
tilt is invertible:
\[
  \tilde\nu = \frac{|A|}{1-\beta}\,\delta_\beta
  = \lambda\,\delta_\beta,
  \qquad \lambda = \frac{|A|}{1-\beta} > 0 .
\]
Therefore $\nu = \lambda\delta_b$ with $b = (\ln\beta)/k < 0$ and
$\lambda > 0$.  The generator $\log W$ is compound Poisson with
deterministic jump size~$b$ and rate~$\lambda$: this is the
log-Poisson distribution.
\end{proof}

\end{document}
